\paperfigure{figure10_topic_stability}{fig:p1f10}{Prespecified capacity and training-seed sensitivities in the original discovery cohort. (A) All nine 40,000-step fits retain their one-to-one assigned match to frozen reference fibroblast-associated topic, without favorable-seed selection. Correlation compares square-root decoder gene probabilities; a missing primary match is unavailable rather than rescued with a rowwise best match. (B) Every reference-topic assignment is shown; gray cells are unmatched or undefined, not low-valued assigned correlations. Capacity K=10 necessarily leaves at least five of the 15 reference topics unmatched. (C) Each point is one existing person's range across three seeds at a fixed K, separately for mean loading and fraction above the fit-specific fibroblast cut. The plotted population is fibroblast-only; complete all-cell counterparts and unavailable counts accompany the tables. Equal-specimen means within patient define this sensitivity experiment, not the primary cell-pooled outcome estimand. (D) Points are person-averaged held-cell log-score differences for a two-component mixture versus a continuous density, on fit-specific transformed topic scores. Positive values favor the mixture under that predictive density comparison. Cell splits and threshold exclusions do not make the full-discovery trained representation independent. Seeds, capacities and points do not add biological patients or validate a discrete state, healing biomarker or amputation endpoint.}
